
\begin{remark}{7.4.106}\label{rem:7-4-106}
定理~\ref{thm:7-4-41}
还有另一种证明，
事实上比上面给出的证明更短。
为完整起见，说明如下。

令
\[
\widehat{\mathfrak h}:
\widehat B(H[1])
\longrightarrow
\widehat B(C[1])
\]
为定理~\ref{thm:5-4-2}
中的同伦等价所诱导的余代数同态。
它的同伦逆诱导余代数同态
\[
\widehat{\mathfrak h}^{-1}:
\widehat B(C[1])
\longrightarrow
\widehat B(H[1]).
\]
复合
\[
\widehat{\mathfrak h}^{-1}
\circ\widehat{\mathfrak h}:
\widehat B(H[1])
\longrightarrow
\widehat B(H[1])
\]
同伦于恒等同态。
由于 $H[1]$ 是典范的，
命题~\ref{prop:5-4-5} 蕴含
$\widehat{\mathfrak h}^{-1}
\circ\widehat{\mathfrak h}$
是同构。
因此，必要时改变
$\widehat{\mathfrak h}^{-1}$，
可以假定
$\widehat{\mathfrak h}^{-1}
\circ\widehat{\mathfrak h}$
就是恒等同态。

注意，
$\widehat{\mathfrak h}^{-1}$
与 $\widehat{\mathfrak h}$
都是余代数同态，
同时也是上链同态。
于是得到
\[
\widehat{\mathfrak H}:
\operatorname{Der}B(H[1])
\longrightarrow
\operatorname{Der}B(C[1]),
\]
其定义为
\[
\widehat{\mathfrak H}(Y)
=
\widehat{\mathfrak h}
\circ Y
\circ\widehat{\mathfrak h}^{-1}.
\]
由于
$\widehat{\mathfrak h}^{-1}
\circ\widehat{\mathfrak h}$
是恒等同态，
$\widehat{\mathfrak H}$
是微分分次李代数同态。
由引理~\ref{lem:7-4-90}，
$\widehat{\mathfrak H}$
是同伦等价，
这正是定理~\ref{thm:7-4-41}
所断言的结论。

不过要指出，
在本小节前面对
定理~\ref{thm:7-4-41} 的证明中，
我们得到一个显式的滤过
$L_\infty$ 同态
\[
\operatorname{Der}B(C[1])
\longrightarrow
\operatorname{Der}B(H[1]).
\]
我们将在第 7.4.6 小节
使用这个滤过 $L_\infty$ 同态的显式描述。

另一方面，
在本注所述的证明中，
我们借助定理 A3.16
在相反方向得到一个微分分次李代数同态，
因而这个证明并非显式的。
\end{remark}
